% from ega3-1-fr.tex
\oldpage[III]{159}

cohomologique), les \emph{extensions} d'un faisceau structural
$\mathcal O_{\mathfrak X}$ (pour un préschéma usuel ou un préschéma formel) par
un Idéal de carré nul, autrement dit les $\mathcal O_{\mathfrak X}$-Algèbres
$\mathcal A$ telles que $\mathcal O_{\mathfrak X}$ soit isomorphe à
$\mathcal A/\mathcal J$, où $\mathcal J$ est un Idéal de carré nul de
$\mathcal A$.
\end{remark}

\subsection{Une décomposition de certains schémas}
\label{subsection:III.5.5-fr}

\begin{proposition}[5.5.1]
\label{III.5.5.1-fr}
Soient $A$ un anneau adique noethérien, $\mathfrak J$ un idéal de définition
de $A$, $Y=\operatorname{Spec}(A)$. Soit $f:X\to Y$ un morphisme séparé de
type fini~; on pose $Y_0=\operatorname{Spec}(A/\mathfrak J)$,
$X_0=X\times_Y Y_0=f^{-1}(Y_0)$. Soit $Z_0$ une partie ouverte de $X_0$,
propre sur $Y_0$~; il existe alors dans $X$ une partie ouverte et fermée $Z$,
propre sur $Y$ et telle que $Z\cap X_0=Z_0$.
\end{proposition}

\begin{proof}
Par hypothèse, il y a un ouvert $T$ de $X$ tel que $T\cap X_0=Z_0$~; soit
$\widehat T$ le complété le long de $Z_0$ du schéma induit par $X$ sur
l'ouvert $T$~; le support de $\mathcal O_{\widehat T}$ étant $Z_0$, qui est
propre sur $Y_0$, $\widehat T$ est propre sur
$\widehat Y=\operatorname{Spf}(A)$ (3.4.1). Il résulte de (5.1.8) qu'il existe
un sous-schéma fermé $Z$ de $T$ propre sur $Y$, tel que, si $i:Z\to T$ est
l'injection canonique, $\widehat i:\widehat Z\to\widehat T$ soit un
isomorphisme ($\widehat Z$ étant le complété de $Z$ le long de $Z_0$). On en
conclut (4.6.8) qu'il existe dans $T$ un voisinage ouvert $V$ de $Z_0$ tel que
la restriction $i^{-1}(V)\to V$ de $i$ soit un isomorphisme. Mais
$i^{-1}(V)$ est un voisinage de $Z_0$ dans $Z$, donc est nécessairement
identique à $Z$ (5.1.3.1). On en conclut que $Z$ est ouvert dans $T$, donc
dans $X$, ce qui achève la démonstration.
\end{proof}

\begin{corollary}[5.5.2]
\label{III.5.5.2-fr}
Si $X_0$ est propre sur $Y_0$, $X$ est réunion de deux parties ouvertes
disjointes $Z$ et $Z'$ telles que $Z$ soit propre sur $Y$ et contienne $X_0$~;
en outre, toute partie fermée $P$ de $X$, propre sur $Y$, est contenue dans
$Z$.
\end{corollary}

\begin{proof}
La dernière assertion résulte de ce que $P\cap Z'$ étant fermé dans $P$ est
propre sur $Y$~; si $P\cap Z'$ n'était pas vide, $f(P\cap Z')$ serait fermé
non vide dans $Y$, donc rencontrerait $Y_0$ (5.1.3.1), ce qui contredit la
définition de $Z$.
\end{proof}

\begin{flushright}
\emph{(A suivre.)}
\end{flushright}

% La page imprimée 504 est blanche.
\clearpage
\thispagestyle{empty}
\null
\clearpage
\thispagestyle{plain}

